\begin{statementbox}
	\textbf{\textcolor{CStatement}{条件}}
	\begin{description}[style=nextline, leftmargin=2.8em, labelsep=0.5em, itemsep=0.35em]
		\item[\textbf{\textcolor{CStatement}{条件 1（两定点 A,B）：}}]
			给定平面上两定点 $A(x_1,y_1)$，$B(x_2,y_2)$ 和动点 $C(x,y)$。
	\end{description}
	\textbf{\textcolor{CStatement}{结论}}
	\begin{description}[style=nextline, leftmargin=2.8em, labelsep=0.5em, itemsep=0.45em]
		\item[\textbf{\textcolor{CStatement}{结论一（以AB为直径的圆）：}}]
			\textbf{\textcolor{CStatement}{直观描述：}}
			直角在点 $C$ 时，$C$ 在以 $AB$ 为直径的圆上（$A,B$ 除外）。
			\[
				\overrightarrow{CA}\cdot\overrightarrow{CB}=0,
				\qquad
				C\neq A,B.
			\]
		\item[\textbf{\textcolor{CStatement}{结论二（过A的垂线）：}}]
			\textbf{\textcolor{CStatement}{直观描述：}}
			直角在点 $A$ 时，$C$ 在过 $A$ 且垂直于 $AB$ 的直线上（$A$ 除外）。
			\[
				\overrightarrow{AB}\cdot\overrightarrow{AC}=0,
				\qquad
				C\neq A.
			\]
		\item[\textbf{\textcolor{CStatement}{结论三（过B的垂线）：}}]
			\textbf{\textcolor{CStatement}{直观描述：}}
			直角在点 $B$ 时，$C$ 在过 $B$ 且垂直于 $AB$ 的直线上（$B$ 除外）。
			\[
				\overrightarrow{BA}\cdot\overrightarrow{BC}=0,
				\qquad
				C\neq B.
			\]
		\item[\textbf{\textcolor{CStatement}{常见变形（圆坐标方程）：}}]
			\textbf{\textcolor{CStatement}{直观描述：}}
			当直角在 $C$ 时，点积展开得直径圆方程，便于联立求解。
			\[
				(x-x_1)(x-x_2)+(y-y_1)(y-y_2)=0.
			\]
	\end{description}

\end{statementbox}
